%!TeX program = pdflatex
% Tailored CV for the 2027--2028 Bloomberg Data Science Ph.D. Fellowship.
% Kept separate from cv-reid.tex, which the GitHub Action publishes publicly.
\documentclass[11pt,letterpaper]{article}
\usepackage[T1]{fontenc}
\usepackage[USenglish]{babel}
\usepackage{palatino}
\usepackage{microtype}
\usepackage[margin=0.85in]{geometry}
\usepackage{enumitem}
\usepackage{titlesec}
\usepackage{fancyhdr}
\usepackage{needspace}
\usepackage[hidelinks]{hyperref}
\urlstyle{same}
\setlength{\parindent}{0pt}
\setlength{\parskip}{2pt}
\titleformat{\section}{\large\bfseries}{}{0pt}{}[\titlerule]
\titlespacing*{\section}{0pt}{9pt}{5pt}
\titlespacing*{\subsection}{0pt}{9pt}{5pt}
\setlist[itemize]{leftmargin=1.1em,itemsep=2pt,parsep=0pt,topsep=2pt}
\newlist{pubs}{enumerate}{1}
\setlist[pubs]{leftmargin=1.7em,label={[\arabic*]},itemsep=5pt,parsep=0pt,topsep=3pt}
\newcommand{\entry}[2]{\Needspace{4\baselineskip}\textbf{#1}\hfill #2\par}
\newcommand{\me}{\textbf{Chen, Y.}}
\newcommand{\earlier}[1]{\\{\small Earlier version: #1}}
\pagestyle{fancy}
\fancyhf{}
\fancyhead[L]{\small Yi Chen}
\fancyhead[R]{\small Curriculum Vitae}
\fancyfoot[L]{\small September 2026}
\fancyfoot[R]{\small\thepage}
\renewcommand{\headrulewidth}{0pt}
\setlength{\headheight}{14pt}
\fancypagestyle{firstpage}{
  \fancyhf{}
  \fancyfoot[L]{\small September 2026}
  \fancyfoot[R]{\small\thepage}
  \renewcommand{\headrulewidth}{0pt}
}
\hypersetup{
  pdfauthor={Yi Chen},
  pdftitle={Yi Chen: Curriculum Vitae},
  pdfkeywords={trustworthy AI, conformal prediction, factuality, retrieval-augmented generation, clarifying questions, preference learning}
}

\begin{document}
\thispagestyle{firstpage}
{\Huge\bfseries Yi Chen}\par
\vspace{5pt}
Ph.D. Candidate in Electrical and Computer Engineering, University of Wisconsin--Madison\\
Expected graduation: 2028\enspace$\vert$\enspace Advisor: \href{https://ramyakv.github.io/}{Ramya Korlakai Vinayak}\\
\href{mailto:yi.chen@wisc.edu}{yi.chen@wisc.edu}\enspace$\vert$\enspace
\href{https://www.deepneural.network/}{deepneural.network}\enspace$\vert$\enspace
\href{https://github.com/kitkatdafu}{github.com/kitkatdafu}\enspace$\vert$\enspace
\href{https://scholar.google.com/citations?user=QoO6pMEAAAAJ}{Google Scholar}

\section*{Research Interests}
\textbf{Trustworthy language models for document question answering and summarization:} conformal factuality and uncertainty quantification for retrieval-augmented generation (RAG) under distribution shift; efficient claim verification; eliciting user intent with few, well-designed questions; learning from heterogeneous human preferences.

\section*{Education}
\textbf{University of Wisconsin--Madison}
\begin{itemize}[label={},leftmargin=0pt,itemsep=2pt]
\item Ph.D., Electrical and Computer Engineering \hfill 2023--present (expected 2028)
\item M.S., Electrical and Computer Engineering \hfill 2022--2025
\item B.A., Computer Sciences, \textit{Honors in the Major} \hfill 2019--2022
\item B.A., Mathematics \hfill 2019--2022
\end{itemize}

\section*{Publications}
\textit{$^\star$ denotes equal contribution.}

\subsection*{Preprint}
\begin{pubs}
\item \me, Chen, D., Chikodikar, S. M., Yin, C. H., and Vinayak, R. K.
\href{https://arxiv.org/abs/2603.16817}{``Is Conformal Factuality for RAG-based LLMs Robust? Novel Metrics and Systematic Insights.''}
\textit{arXiv:2603.16817}, 2026.
\end{pubs}

\subsection*{Conference Publications}
\begin{pubs}[resume]
\item Tatli, G., \me, Mason, B., Nowak, R. D., and Vinayak, R. K.
\href{https://proceedings.mlr.press/v286/tatli25a.html}{``Metric Learning in an RKHS.''}
\textit{Conference on Uncertainty in Artificial Intelligence (UAI)}, 2025.

\item Vishwakarma, H.$^\star$, \me$^\star$, Namburi GNVV, S. S. S., Tay, S. J., Vinayak, R. K., and Sala, F.
\href{https://proceedings.mlr.press/v267/vishwakarma25a.html}{``Rethinking Confidence Scores and Thresholds in Pseudolabeling-based SSL.''}
\textit{International Conference on Machine Learning (ICML)}, 2025.

\item \me\ and Vinayak, R. K.
\href{https://doi.org/10.1145/3696410.3714587}{``Query Design for Crowdsourced Clustering: Effect of Cognitive Overload and Contextual Bias.''}
\textit{The ACM Web Conference (WWW)}, 2025. \textbf{Oral presentation.}
\earlier{ICML Workshop on Models of Human Feedback for AI Alignment, 2024.}

\item Chen, D., \me, Rege, A., Wang, Z., and Vinayak, R. K.
\href{https://openreview.net/forum?id=1kFDrYCuSu}{``PAL: Pluralistic Alignment Framework for Learning from Heterogeneous Preferences.''}
\textit{International Conference on Learning Representations (ICLR)}, 2025.

\item Vishwakarma, H., \me, Tay, S. J., Namburi, S. S. S., Sala, F., and Vinayak, R. K.
\href{https://proceedings.neurips.cc/paper_files/paper/2024/hash/1d051fb631f104cb2a621451f37676b9-Abstract-Conference.html}{``Pearls from Pebbles: Improved Confidence Functions for Auto-labeling.''}
\textit{Advances in Neural Information Processing Systems (NeurIPS)}, 2024.

\item Tatli, G., \me, and Vinayak, R. K.
\href{https://proceedings.mlr.press/v238/tatli24a.html}{``Learning Populations of Preferences via Pairwise Comparison Queries.''}
\textit{International Conference on Artificial Intelligence and Statistics (AISTATS)}, 2024.
\earlier{ICML MFPL Workshop, 2023.}

\item \me, Vinayak, R. K., and Hassibi, B.
\href{https://doi.org/10.1609/hcomp.v11i1.27545}{``Crowdsourced Clustering via Active Querying: Practical Algorithm with Theoretical Guarantees.''}
\textit{AAAI Conference on Human Computation and Crowdsourcing (HCOMP)}, 2023.
\earlier{ICML AI \& HCI Workshop, 2023.}
\end{pubs}

\subsection*{Workshop Publications}
\begin{pubs}[resume]
\item \me, Yin, C. H., Chikodikar, S. M., and Vinayak, R. K.
\href{https://openreview.net/forum?id=BYVo4s8AcE}{``On the Scoring Functions for RAG-based Conformal Factuality.''}
\textit{ICML Workshop on Reliable and Responsible Foundation Models}, 2025.

\item Vishwakarma, H., \me, Namburi GNVV, S. S. S., Tay, S. J., Vinayak, R. K., and Sala, F.
\href{https://openreview.net/forum?id=2Nd3qXcAnJ}{``PabLO: Improving Semi-Supervised Learning with Pseudolabeling Optimization.''}
\textit{ICML Workshop on Data-Centric Machine Learning Research}, 2024.

\item Tatli, G., \me, and Vinayak, R. K.
``Learning Preference Distributions From Pairwise Comparisons.''
\textit{International Workshop on Computational Social Choice (COMSOC)}, 2023.
\end{pubs}

\section*{Research Experience}
\entry{Conformal Factuality for Retrieval-Augmented Generation}{2025--present}
\begin{itemize}
\item Led a study of conformal factuality for RAG across three benchmarks and multiple model families; introduced informativeness-aware metrics and diagnosed vacuous outputs, sensitivity to retrieval shift, and sensitivity to distractors.
\item Found that lightweight entailment verifiers match or outperform LLM confidence scorers with over $100\times$ fewer FLOPs (preprint, 2026; ICML workshop, 2025).
\item Ongoing: claim verification with masked-diffusion language models against trained checkers and token-probability baselines.
\end{itemize}
\entry{Human Feedback and Preference Learning}{2023--2025}
\begin{itemize}
\item Designed large crowdsourced studies showing how question design trades cognitive overload against contextual bias (WWW 2025, oral); developed active querying with recovery guarantees (HCOMP 2023).
\item Co-developed PAL, which learns shared structure and individual differences from heterogeneous preferences (ICLR 2025), and theory for learning preference distributions from pairwise comparisons (AISTATS 2024; UAI 2025).
\item Co-developed confidence functions and jointly designed scores and thresholds for auto-labeling and pseudolabeling under error constraints (NeurIPS 2024; ICML 2025).
\end{itemize}

\section*{Open-Source Software and Data}
\begin{itemize}
\item \textbf{PAL} (ICLR 2025): implementation and datasets, Apache-2.0.\\
\href{https://github.com/RamyaLab/pluralistic-alignment}{\nolinkurl{github.com/RamyaLab/pluralistic-alignment}}
\item \textbf{Query design} (WWW 2025): code and crowdsourced data, archived on Zenodo.\\
\href{https://github.com/kitkatdafu/cognitive-overload-contextual-bias}{\nolinkurl{github.com/kitkatdafu/cognitive-overload-contextual-bias}}
\item \textbf{Populations of preferences} (AISTATS 2024): code and data.\\
\href{https://github.com/ramyakv/Learning-Population-of-Preferences-AISTATS2024}{\nolinkurl{github.com/ramyakv/Learning-Population-of-Preferences-AISTATS2024}}
\item \textbf{Active crowdsourced clustering} (HCOMP 2023): code and data.\\
\href{https://github.com/kitkatdafu/crowd-active-clustering}{\nolinkurl{github.com/kitkatdafu/crowd-active-clustering}}
\end{itemize}

\Needspace{6\baselineskip}
\section*{Industry Experience}
\entry{Meta --- Core Ads Growth Engineering}{2026}
Software Engineer Intern, Machine Learning (Summer 2026)\\
Part-time Student Researcher (Fall 2026)
\begin{itemize}
\item Developed a multi-turn LLM agent for generating and editing ad creatives from advertiser briefs, together with a vision-language-model-as-judge evaluation suite.
\end{itemize}

\section*{Teaching and Mentoring}
\textbf{University of Wisconsin--Madison}\\
Teaching Assistant, ECE 539: Introduction to Artificial Neural Networks\\
Summer 2025; Fall 2025, 2026
\begin{itemize}
\item Mentored undergraduate researchers on research projects and scientific writing.
\end{itemize}

\section*{Professional Service}
\textbf{Conference reviewer:} NeurIPS 2026; COLM 2026; UAI 2026; ICLR 2025, 2026; WWW 2025.\\
\textbf{Workshop reviewer:} NeurIPS 2024; ICML 2024.

\section*{Awards and Scholarships}
HCOMP/CI Travel Award, 2023; UW--Madison Conference Presentation Award; UW--Madison Undergraduate Scholarship for Summer Study; James Graaskamp Scholarship.

\section*{Technical Skills}
\textbf{Research methods:} Conformal prediction and uncertainty quantification; factuality evaluation for RAG; LLM-judge evaluation; human-subject studies and preference learning.\\
\textbf{Machine learning:} PyTorch, Hugging Face Transformers, JAX, TensorFlow, scikit-learn; supervised fine-tuning, model serving.\\
\textbf{Programming:} Python (primary), C/C++, CUDA, Java, JavaScript, Rust, SQL.\\
\textbf{Systems:} Docker, Git, AWS (EC2, S3, Lambda); React, Node.js and MongoDB for crowdsourcing interfaces.

\end{document}
